\section{Build a large index with bounded work buffers}
\label{sec:index-construction}

A simple MPS builder could collect every $(hash,position,strand)$ record and
sort the complete vector. Its memory use would then grow with the number of
target minimizers. BioSeqDB instead sorts smaller \emph{runs}: ordered groups
of records that can be merged without holding their complete contents in
memory.

The source scan must first handle boundaries correctly. A \emph{halo} is
the extra context read on either side of a core. For the current grouped
builder, a halo of $k+w-2$ symbols covers the span needed around a minimizer
window. An occurrence belongs to the core containing its last-symbol
coordinate. Thus overlap supplies context without giving the same occurrence
to two owners.

\begin{figure}[htbp]
\centering
\begin{tikzpicture}[x=1cm,y=1cm,font=\footnotesize]
  \node[anchor=west,font=\bfseries] at (0,5.1) {Neighboring construction cores};
  \draw[fill=indexBlue!12,draw=indexBlue] (1,3.25) rectangle (6,3.85);
  \draw[fill=indexGreen!12,draw=indexGreen] (6,3.25) rectangle (11,3.85);
  \node at (3.5,3.55) {core 0}; \node at (8.5,3.55) {core 1};
  \draw[dashed,black!55] (6,.25)--(6,4.15);
  \node[above] at (6,4.15) {boundary $b$};
  \draw[fill=indexBlue!8,draw=indexBlue] (1,2.15) rectangle (7.2,2.62);
  \draw[pattern=north east lines,pattern color=indexBlue,draw=indexBlue]
    (6,2.15) rectangle (7.2,2.62);
  \draw[fill=indexGreen!8,draw=indexGreen] (4.8,1.05) rectangle (11,1.52);
  \draw[pattern=north east lines,pattern color=indexGreen,draw=indexGreen]
    (4.8,1.05) rectangle (6,1.52);
  \node[anchor=east] at (.8,2.38) {read 0};
  \node[anchor=east] at (.8,1.28) {read 1};
  \node[anchor=west] at (7.45,2.38) {right halo};
  \node[anchor=east] at (4.55,1.28) {left halo};
  \draw[very thick,indexOrange] (5.25,.5)--(6.25,.5);
  \fill[indexOrange] (6.25,.5) circle[radius=2.5pt];
  \node[anchor=west] at (6.5,.5) {seed ending in core 1: emitted by core 1};
\end{tikzpicture}
\caption{The readers overlap, but their output ownership does not. The
orange seed crosses the boundary; its last position belongs to core 1.
The drawing shows only the internal halos, with lengths exaggerated.}
\label{fig:index-build-halos}
\end{figure}

\noindent\begin{minipage}{\linewidth}
MPS turns the resulting global records into a sequence-wide order:
\begin{enumerate}
  \item Fill one work buffer, sort it by hash, then position and strand, and
    remove duplicate records.
  \item Write the sorted run to PostgreSQL temporary-file storage, then
    reuse the work buffer.
  \item Merge a bounded number of runs at a time. Reading a small prefix
    from each input is enough to choose the next output record.
  \item Continue merging until one ordered stream remains. Remove
    duplicates across runs as well.
  \item Read that stream to count each hash's occurrences, then replay it
    to encode controls and positions with the correct counts.
\end{enumerate}
\end{minipage}

\begin{figure}[htbp]
\centering
\begin{tikzpicture}[x=1cm,y=1cm,font=\footnotesize]
  \node[anchor=west,font=\bfseries] at (0,5.8) {Small sorted runs};
  \foreach \x/\name/\items in {1.4/0/{a,c,f},4.5/1/{a,b,e},7.6/2/{b,d,f},10.7/3/{a,c,e}} {
    \node[index blue,minimum width=2.5cm] (r\name) at (\x,4.8) {run \name\\$\items$};
  }
  \node[index orange,minimum width=3.6cm] (m0) at (3,3.0) {$a,a,b,c,e,f$};
  \node[index orange,minimum width=3.6cm] (m1) at (9,3.0) {$a,b,c,d,e,f$};
  \draw[index arrow] (r0.south)--(m0.north west);
  \draw[index arrow] (r1.south)--(m0.north east);
  \draw[index arrow] (r2.south)--(m1.north west);
  \draw[index arrow] (r3.south)--(m1.north east);
  \node[anchor=west,align=left] at (0,1.7)
    {Only small read buffers\\for each open input};
  \node[index green,minimum width=6.8cm] (final) at (7.6,1.2)
    {$a,a,a,b,b,c,c,d,e,e,f,f$};
  \draw[index arrow] (m0.south)--(final.north west);
  \draw[index arrow] (m1.south)--(final.north east);
  \node[anchor=west,text=black!65] at (0,.1)
    {Two runs merge at a time. Only hashes are shown; positions are omitted.};
\end{tikzpicture}
\caption{External merging bounds the number of simultaneously open input
runs. Equal letters represent equal hashes with different positions, so
their occurrences are preserved. Only duplicate full records are removed.}
\label{fig:index-mps-merge}
\end{figure}

The spill implementation uses PostgreSQL's \texttt{BufFile} temporary-file
interface, rather than giving each run an unmanaged permanent file. The
number of merge inputs is limited independently of the total run count.
One unusually frequent hash can also have a posting list larger than a work
buffer. Its positions are streamed through bounded buffers, so one long list
does not require one equally large vector.
Temporary disk use can still grow with the input and the merge passes. A
memory bound does not imply constant disk space or constant build time.

Construction writes the complete payload before its final navigation
metadata becomes usable. In the sequence-wide stream layout, directory
entries may only refer to finished controls and postings. Publication of the
new generation is a separate final step, described in
Section~\ref{sec:index-generations}.
